\subsubsection{Bidirectional electronic transport and arithmetic record}
\label{ii:transporte-electronico}

The central representation uses a cell shared by the two
APP evaluations. Its coordinates belong to \(D_9=\{1,\ldots,9\}\),
and the inversion \((r,c)\mapsto(10-r,10-c)\) has the unique fixed point
\((5,5)\). The visible state is written \((r,c,h,v,n)\), with
\(h,v\in\{-1,1\}\) and integer chronological index \(n\geq0\).
The tritic selector repeats \((+1,-1,0)\). On the additive sheet,
\(r+c\) is read and the column is translated in direction \(h\); on the
multiplicative sheet, \(rc\) is read and the row is translated in direction
\(v\). The threshold preserves position. Both orientations are
reversed upon completing each interval of twenty-seven updates.

To express reading and carry with the same positive convention,
define
\[
 [z]^+_9=1+((z-1)\bmod9),\qquad
 q_9(z)=\frac{z-[z]^+_9}{9}.
\]
Evaluation precedes displacement. If \(a_n\) is the integer
read, the event preserves \((a_n,[a_n]^+_9,q_9(a_n))\), the sheet,
previous position, orientation, and boundary crossing. The
tritic symbol is \([a_n]^+_9\bmod3\). The spatial lift
retains, for example, the integer
\(q_9(c+h)\) when \(c\) is replaced by \([c+h]^+_9\).

The shared-cell representation and the two-cursor
representation of \eqref{ii:tpk-fobs-posiciones} have different visible
states. Their records are related to the TPK state through
their respective representation maps. For the central trajectory,
the full update on the shared cell producing the electronic word
is used here.

\begin{proposition}[Invertible update with preserved quotients]
\label{ii:electron-inversa-cocientes}
Add to the visible state two integer records \(Q_+,Q_\times\)
and two winding numbers \(W_r,W_c\). Each active reading
increments its sheet's record by \(q_9(a_n)\); each displacement
increments the corresponding winding by its boundary crossing.
This update has an explicit inverse on its image.
\end{proposition}
\begin{proof}
The later index determines the preceding mode and whether an orientation
reversal has just occurred. This reversal is undone first.
If the mode was additive, the previous column is recovered as
\([c-h]^+_9\); if multiplicative, the row is recovered as
\([r-v]^+_9\). The neutral mode preserves both coordinates. The
recovered cell again determines the integer read, its remainder,
quotient, and boundary crossing. Subtracting these increments from
\(Q_+,Q_\times,W_r,W_c\) and reversing the chronological index
recovers the preceding state exactly. Composing these
inverses recovers an entire finite history. Event concatenation
also preserves the order of readings.
\end{proof}

The records \(Q_+,Q_\times\) are an explicit representation by sums
of the quotient history. The construction is integrated into the memory
fiber alongside its other coordinates. Its composition law,
for concatenable histories \(u,v\), is
\[
 Q_s(uv)=Q_s(u)+Q_s(v),\qquad s\in\{+,\times\},
\]
where the second history begins at the first one's final state.
Algebraic inversion of an update subtracts its increments;
subsequently traversing a spatial return path records
new evaluations. This distinction simultaneously preserves
transport reversibility and history duration.

\begin{proposition}[Central return and exact accumulation]
\label{ii:electron-memoria-retorno}
Starting from \((r,c,h,v,n)=(5,5,1,1,0)\) and zero initial records, the
spatial projection and its orientations return after fifty-four
updates. Over that interval,
\[
 (\Delta W_r,\Delta W_c)=(0,0),\qquad
 (\Delta Q_+,\Delta Q_\times)=(10,46).
\]
After \(54N\) updates, for every integer \(N\geq0\),
one has \((Q_+,Q_\times)=(10N,46N)\), while the spatial
coordinates and their orientations return. The integer index takes the value
\(54N\); the dodecaphasic calendar modulo \(108\) returns for
even \(N\).
\end{proposition}
\begin{proof}
Each pair of active modes displaces each coordinate once. In
twenty-seven updates, both complete a turn of nine
positions and return to the central diagonal; \(h,v\) are then
reversed. The next block traverses the opposite turn. Spatial
crossings cancel between the two blocks.

In each block, the additive readings run through \(2k\),
\(1\leq k\leq9\), whose positive quotient sums to five. The
multiplicative readings of the positive block run through
\(k[k+1]^+_9\). Their integer sum and remainder sum are
\[
 \sum_{k=1}^9 k[k+1]^+_9=249,\qquad
 \sum_{k=1}^9[k[k+1]^+_9]^+_9=42.
\]
The second identity follows from the list
\((2,6,3,2,3,6,2,9,9)\). By Euclidean division,
the quotient sum is \((249-42)/9=23\). Substitution
\(k\mapsto[k-1]^+_9\) shows that the negative block contains
the same multiset of products. The two blocks therefore produce
ten and forty-six units in the respective records.
The spatial projection and its orientations return to their initial
values. The activation and reversal rules repeat upon
shifting the index by \(54\), and the increments do not depend on
the accumulated records. The preceding additive law
proves the formula for every \(N\) by induction.
\end{proof}

The tritic reading and window orientation provide, on
the same trajectory,
\[
 \gamma_e=\bigl((100000220)^3(120200000)^3\bigr)^2,\qquad
 \tau_e=(+,+,+,-,-,-,+,+,+,-,-,-).
\]
Powers represent concatenations. The discrepancy
and window readers developed in
\eqref{eq:el-kd} and \eqref{eq:el-komega} both compute
\((80,54,6)\). The evaluation is obtained from the trajectory before
incorporating \(\alpha\), the action section, or a mass.
Quotient memory preserves information additional to that triple:
in one hundred eight updates it records \((20,92)\), with
zero net spatial winding. Each representation thus retains the
type of the observable producing it.
